\scanpage{17}
Éliminons de $=$ le cas $g=0$, $\nu\leq2$. Ici $g=0=\sum g_\alpha+h_1$ implique que les $g_\alpha$ et $h_1$ sont nuls : tous les sommets sont de genre $0$, le graphe est un arbre. Il ne peut être réduit à un seul sommet, donc il a au moins deux bouts, dont les \unsure{poids} marqués doivent être $\geq2$, donc \deleted{\unsure{le}} on a $\nu\geq2+2=4$, absurde.

Soit $G$ une \deleted{\inserted{MD-}graphe}\inserted{maquette}. On dit qu'une \deleted{\inserted{MD-}} courbe \deleted{\unsure{standard|standart}}\inserted{\unsure{standard|standart}} $C$ sur un corps alg. clos est de \uline{type $G$}, si sa maquette est iso à $G$, on dit qu'elle est \inserted{$G$-}\uline{épinglée} si on se donne un iso entre sa maquette et $G$ (c'est donc une structure suppl. sur $\unsure{C|X}$, appelée $G$-épinglage).

Soit $(\widehat X,I)$ une courbe \unsure{standard|standart} sur une base $S$ quelconque, on dit qu'elle est de type $G$ si ses fibres géom. sont de type $G$. Alors les \deleted{\unsure{``types'' des}}\inserted{maquettes} des
\mspage{16}
